Bagian ini menyajikan rencana jadwal penelitian beserta tahapan dan estimasi waktu pelaksanaan, mengikuti rencana yang telah dipresentasikan dan disetujui pembimbing. Jadwal dapat disesuaikan mengikuti kalender akademik. Matriks kegiatan ditunjukkan pada Tabel \ref{tab:jadwal_penelitian}.

\begin{table}[!h]
    \centering
    \caption{Jadwal Kegiatan Penelitian}
    \label{tab:jadwal_penelitian}
    \resizebox{\linewidth}{!}{
    \begin{tabular}{clcccccccccc}
    \toprule
        No & Kegiatan & Mei & Jun & Jul & Agu & Sep & Okt & Nov & Des & Jan & Feb \\
    \midrule
        1 & Identifikasi masalah \& studi literatur & X & X & & & & & & & & \\
        2 & Penyusunan proposal tesis & & X & X & X & & & & & & \\
        3 & Pengumpulan \& penyusunan dataset & & & & X & X & X & & & & \\
        4 & \textit{Fine-tuning} \textit{Local} SLM (LoRA) & & & & & & X & X & & & \\
        5 & Implementasi RAG \& \textit{re-ranking} & & & & & & & X & X & X & \\
        6 & Pengujian \& evaluasi sistem & & & & & & & & X & X & \\
        7 & Analisis hasil penelitian & & & & & & & & & X & X \\
        8 & Penulisan tesis \& revisi & & & & & & & & X & X & X \\
    \bottomrule
    \end{tabular}
    }
\end{table}

Keterangan: tanda ``X'' menunjukkan periode pelaksanaan tiap kegiatan. % TODO (Aditya): sesuaikan tahun dan bulan mengikuti kalender akademik serta arahan pembimbing.
